% =============================================================================
\section{Protocolo y seguridad}
% =============================================================================

\begin{frame}[fragile]{El protocolo: texto, una orden por línea}
\begin{columns}[T,onlytextwidth]
\begin{column}{0.47\textwidth}
\footnotesize
\textbf{Del Pi al ESP32}
\begin{verbatim}
M,<v_lin>,<v_ang>    enteros de -100 a 100
B,<patron>           0 apagado, 1 hallazgo,
                     2 alarma
\end{verbatim}
\textbf{Del ESP32 al Pi} (20 veces por segundo)
\begin{verbatim}
D,<dist_cm>,<enc_izq>,<enc_der>
# ...                diagnostico: se descarta
\end{verbatim}
\textbf{Una sesión real, resumida}
\begin{verbatim}
<- D,153,0,0         pared a 1,5 m
-> M,0,0             mirando
-> M,0,35            un paso de giro
-> M,0,0             mirando
-> M,55,12           hacia el oso
-> M,0,0             llego
-> B,1               hallazgo: pitido
\end{verbatim}
\end{column}
\begin{column}{0.51\textwidth}
\footnotesize
\begin{block}{Por qué en texto}
  Se lee con cualquier monitor serie y se depura a ojo. A 115200 baudios, una
  línea de 10 caracteres tarda menos de 1 ms: el protocolo no es el cuello de
  nada.
\end{block}
\begin{block}{Las reglas chicas que importan}
  \begin{itemize}
    \item \verb|v_ang| positivo es girar a la \textbf{derecha}.
    \item \verb|dist_cm = -1| es ``sin lectura'', \textbf{no} ``distancia
          cero''. Tratarlo como cero frenaría el robot de entrada.
    \item Los encoders se mandan en cero en vez de omitirse: el Pi lee siempre
          el mismo formato.
    \item Las líneas que empiezan con \verb|#| son para humanos; el Pi las
          cuenta y las tira.
    \item Para parar, exactamente \verb|M,0,0|.
  \end{itemize}
\end{block}
\end{column}
\end{columns}
\end{frame}

\begin{frame}{Las dos redes de seguridad, en el tiempo}
  \centering
  \begin{tikzpicture}[font=\scriptsize, x=1.65cm, y=0.55cm]
    \draw[->, gris, line width=0.8pt] (0,0) -- (7.0,0) node[right]{tiempo};
    \foreach \x/\t in {0/0, 2/0{,}40 s, 2.5/0{,}50 s, 3.5/0{,}70 s} {
      \draw[gris] (\x,0.15) -- (\x,-0.15) node[below]{\t};
    }
    \node[above, align=center, text width=2.6cm] at (0,0.3) {el último comando\\válido};
    \node[above, align=center, text width=3.0cm, azul] at (2,0.3) {\textbf{vigencia}: el latido\\manda \cod{M,0,0}};
    \node[above, align=center, text width=3.2cm, naranja] at (3.6,2.3) {\textbf{watchdog}: el ESP32\\frena activo};
    \draw[naranja] (2.5,0.15) -- (2.9,2.3);
    \node[above, align=center, text width=2.6cm, gris] at (5.4,0.3) {suelta a\\rueda libre};
    \draw[gris] (3.5,0.15) -- (4.7,0.9);
  \end{tikzpicture}

  \vspace{0.4em}
  \relax{\footnotesize
  \begin{columns}[T,onlytextwidth]
    \begin{column}{0.32\textwidth}
      \begin{block}{Primera red: la vigencia}
        Actúa si el \textbf{programa sigue vivo} pero el lazo de visión dejó
        de decidir (un frame colgado, la cámara trabada). La pone el Pi.
      \end{block}
    \end{column}
    \begin{column}{0.32\textwidth}
      \begin{block}{Segunda red: el watchdog}
        Actúa si el Pi \textbf{dejó de hablar}: programa muerto, cable
        desconectado, SSH cortado. La pone el ESP32 y no depende de nadie.
      \end{block}
    \end{column}
    \begin{column}{0.32\textwidth}
      \begin{alertblock}{Tercera: la mano}
        El conector de 12 V. Es la única que sirve si el robot hace algo
        \emph{coherente pero equivocado}, como avanzar contra una pared.
      \end{alertblock}
    \end{column}
  \end{columns}}

  \vspace{0.2em}
  {\footnotesize Las dos primeras se probaron en hardware el 22 de septiembre,
  con el robot elevado: dejando el programa vivo pero callado, y matándolo con
  \cod{kill -9}.}
\end{frame}

\begin{frame}{Dos trampas de instrumentación que ya costaron tiempo}
  \begin{columns}[T,onlytextwidth]
    \begin{column}{0.485\textwidth}
      \begin{alertblock}{Abrir el puerto reinicia el ESP32}
        {\small El puente USB (CP2102) baja la línea DTR al abrir el puerto, y
        eso reinicia el ESP32, que arranca con los motores frenados.

        La primera prueba del watchdog estaba mal por esto: se mataba el
        programa y se abría el puerto ``para ver'' si el ESP32 frenaba\dots{}
        y abrirlo \textbf{ya lo frenaba}. Un instrumento que altera lo que
        mide.}
      \end{alertblock}
    \end{column}
    \begin{column}{0.485\textwidth}
      \begin{alertblock}{Los 2 segundos del arranque}
        {\small Por lo anterior, \cod{Enlace} espera 2 s después de abrir, hasta
        que el ESP32 termina de reiniciarse.

        Si se lanza algo con \cod{timeout N}, el movimiento dura $N - 2$
        segundos. Parece una falla del enlace y no lo es.}
      \end{alertblock}
    \end{column}
  \end{columns}

  \vspace{0.5em}
  \begin{block}{Y una mejora que salió del barrido a pasos}
    {\small \cod{mover()} sólo actualizaba el comando y dejaba que el latido de 10 Hz
    lo mandara: cada cambio tardaba entre 0 y 0,1 s en salir, al azar. Para
    avanzar no importa. Para un pulso de giro de 0,25 s es hasta el 40 \% en
    cada punta: un paso giraba la mitad o el doble que el anterior. Ahora
    \textbf{los cambios de comando se mandan en el acto} y el latido sólo
    sostiene. Como escriben dos hilos, la escritura va con candado.}
  \end{block}
\end{frame}
